% ---------------- เอกสารอ้างอิง ----------------
\chapter*{เอกสารอ้างอิง}
\addcontentsline{toc}{chapter}{เอกสารอ้างอิง}
\markboth{เอกสารอ้างอิง}{เอกสารอ้างอิง}

\begin{enumerate}[leftmargin=2.2em, itemsep=0.75em, label={[\arabic*]}]

\item Jin, Z., Du, X., Xu, Y., Deng, Y., Liu, M., Zhao, Y., \ldots\ \& Yang, H. (2020). Structure of \Mpro\ from SARS-CoV-2 and discovery of its inhibitors. \textit{Nature}, 582(7811), 289--293. \url{https://doi.org/10.1038/s41586-020-2223-y}

\item Gaulton, A., Hersey, A., Nowotka, M., Bento, A. P., Chambers, J., Mendez, D., \ldots\ \& Overington, J. P. (2017). The ChEMBL database in 2017. \textit{Nucleic Acids Research}, 45(D1), D945--D954. \url{https://doi.org/10.1093/nar/gkw1074}

\item Rogers, D., \& Hahn, M. (2010). Extended-connectivity fingerprints. \textit{Journal of Chemical Information and Modeling}, 50(5), 742--754. \url{https://doi.org/10.1021/ci100050t}

\item Lipinski, C. A., Lombardo, F., Dominy, B. W., \& Feeney, P. J. (2001). Experimental and computational approaches to estimate solubility and permeability in drug discovery and development settings. \textit{Advanced Drug Delivery Reviews}, 46(1-3), 3--26. \url{https://doi.org/10.1016/S0169-409X(00)00129-0}

\item Nantasenamat, C., Isarankura-Na-Ayudhya, C., Naenna, T., \& Prachayasittikul, V. (2010). A practical overview of quantitative structure-activity relationship. \textit{EXCLI Journal}, 9, 74--93. \url{https://www.ncbi.nlm.nih.gov/pmc/articles/PMC5084994/}

\item Breiman, L. (2001). Random Forests. \textit{Machine Learning}, 45(1), 5--32. \url{https://doi.org/10.1023/A:1010933404324}

\item Chen, T., \& Guestrin, C. (2016). XGBoost: A scalable tree boosting system. In \textit{Proceedings of the 22nd ACM SIGKDD International Conference on Knowledge Discovery and Data Mining} (pp. 785--794). \url{https://doi.org/10.1145/2939672.2939785}

\item Rodríguez-Pérez, R., \& Bajorath, J. (2020). Interpretation of machine learning predictions using Shapley values: Application to compound activity and property prediction. \textit{Molecular Informatics}, 39(1-2), 1900131. \url{https://doi.org/10.1002/minf.201900131}

\item Bemis, G. W., \& Murcko, M. A. (1996). The properties of known drugs. 1. Molecular frameworks. \textit{Journal of Medicinal Chemistry}, 39(15), 2887--2893. \url{https://doi.org/10.1021/jm9602928}

\item Sheridan, R. P. (2013). Time-split cross-validation as a method for estimating the goodness of prospective prediction. \textit{Journal of Chemical Information and Modeling}, 53(4), 783--790. \url{https://doi.org/10.1021/ci400084k}

\item Lundberg, S. M., \& Lee, S. I. (2017). A unified approach to interpreting model predictions. In \textit{Advances in Neural Information Processing Systems 30} (pp. 4765--4774).

\item Maggiora, G. M. (2006). On outliers and activity cliffs --- Why QSAR often disappoints. \textit{Journal of Chemical Information and Modeling}, 46(4), 1535. \url{https://doi.org/10.1021/ci060117s}

\item Pedregosa, F., Varoquaux, G., Gramfort, A., Michel, V., Thirion, B., Grisel, O., \ldots\ \& Duchesnay, É. (2011). Scikit-learn: Machine learning in Python. \textit{Journal of Machine Learning Research}, 12, 2825--2830.

\item Landrum, G. (2024). \textit{RDKit: Open-source cheminformatics}. \url{https://www.rdkit.org}

\end{enumerate}
